% note éditée à la main (coupe du 17/11/2026) : graphiques.fin ne l'écrase pas
DGFiP, IRCOM et déclarations 2042 des revenus 2024 (répartition) ; Insee, comptes nationaux 2025, et rapport « Financement » annexé au projet de loi d'approbation des comptes de la Sécurité sociale 2025 (agrégats de 2025) ; OpenFisca-France ; calculs de l'auteur (synthese/modele\_purs.py, naive\_ircsg.py).\par Note : réforme naïve (barème de l'IR du foyer et taux proportionnel de 3,99 \%), à recette égale (230 Md\euro{}, base 2025) avant réaction des contribuables. Le 1 \% le plus aisé paierait 2,6 Md\euro{} de plus.
